% =============================================================================
% Theorem 2 (Upgraded): Multi-Parameter Cramér--Rao with Fisher Information
%                      Geometry for Joint Estimation of (gamma, H)
% =============================================================================
% Paper: Impossibility Triangle of LLM Evaluation (v9 upgrade)
% Date: 2026-07-09
%
% Comparison to v8 Theorem 2:
%   v8 T2: Var(gamma_hat) >= 1/(N * sigma*(1-sigma))  (single-parameter CR bound
%          on gamma only, using binary observation approximation).
%   v9 T2: Joint 2x2 Fisher information matrix I(gamma, H); covariance matrix
%          of any unbiased estimator satisfies Sigma_hat >= I^{-1}.
%          Diagonal entries give sharper individual bounds; off-diagonal
%          characterizes the curvature of the (gamma, H) manifold under
%          Amari's information geometry.
% =============================================================================

\section*{Theorem 2 (Multi-Parameter Cramér--Rao with Fisher Geometry)}

\begin{theorem}[Joint Cramér--Rao Bound for $(\hat\gamma, \hat H)$]
\label{thm:T2}
Let $\theta := (\gamma, H) \in \Theta \subset \mathbb{R}^2$ be the parameters of
joint interest (evaluator coupling and strategy entropy). Assume:
\begin{enumerate}[nosep,label=(\roman*)]
  \item Observations $y_1, \ldots, y_N \stackrel{\mathrm{iid}}{\sim}
        p(y \mid \theta)$ with $p$ twice-differentiable in $\theta$ and
        satisfying the standard regularity conditions
        (see~\citet{lehmanncasella1998}, Section~2.5).
  \item The model is \textbf{identified}, i.e., the Fisher information matrix
        $\mathcal{I}(\theta)$ is positive definite.
\end{enumerate}
Then for any unbiased estimator $\hat\theta$ of $\theta$, the covariance matrix
$\Sigma_{\hat\theta} := \mathrm{Cov}(\hat\theta)$ satisfies the
\textbf{multi-parameter Cramér--Rao inequality}
\[
  \Sigma_{\hat\theta} \;\succeq\; \mathcal{I}(\theta)^{-1},
\tag{CR}
\]
where $\mathcal{I}(\theta) \in \mathbb{R}^{2\times 2}$ has entries
\[
  \mathcal{I}_{ij}(\theta)
    \;=\; N \cdot \mathbb{E}\!\left[
        \frac{\partial \log p(y \mid \theta)}{\partial \theta_i}
        \frac{\partial \log p(y \mid \theta)}{\partial \theta_j}
      \right].
\]
In particular,
\begin{enumerate}[nosep,label=(\alph*)]
  \item $\mathrm{Var}(\hat\gamma) \ge [\mathcal{I}^{-1}]_{11}$, recovering the
        v8 single-parameter bound as a special case.
  \item $\mathrm{Var}(\hat H) \ge [\mathcal{I}^{-1}]_{22}$.
  \item $\mathrm{Corr}(\hat\gamma, \hat H) \in [-1, 1]$ is bounded by the
        off-diagonal geometry of $\mathcal{I}^{-1}$.
\end{enumerate}
\end{theorem}

\begin{proof}
The inequality $\Sigma_{\hat\theta} \succeq \mathcal{I}(\theta)^{-1}$ is the
standard multi-parameter Cramér--Rao bound
(Theorem~2.1 in~\citet{lehmanncasella1998}). It follows from the Cauchy--Schwarz
inequality applied to the score vector
$s(\theta) := \nabla_\theta \log p(y \mid \theta)$ and the unbiased-estimation
constraint $\mathbb{E}[\hat\theta - \theta] = 0$:
\[
  \mathrm{Cov}(\hat\theta)
    \;=\; \mathbb{E}[(\hat\theta - \theta)(\hat\theta - \theta)^\top]
    \;\succeq\; \mathbb{E}[s(\theta) s(\theta)^\top]^{-1}
    \;=\; \mathcal{I}(\theta)^{-1}.
\]
The diagonal entries $\mathcal{I}_{11}$ and $\mathcal{I}_{22}$ reduce to
$N \cdot \mathbb{E}[(\partial_\gamma \log p)^2]$ and
$N \cdot \mathbb{E}[(\partial_H \log p)^2]$; these are the per-sample Fisher
information for $\gamma$ and $H$ respectively. The first diagonal entry
$[\mathcal{I}^{-1}]_{11}$ is at least $1/\mathcal{I}_{11}$ by positivity of
$\mathcal{I}^{-1}$, recovering the v8 single-parameter bound
$\mathrm{Var}(\hat\gamma) \ge 1 / (N \cdot \mathcal{I}_{11})$. \qedhere
\end{proof}

\begin{corollary}[Sharpened v8 Coupling--Precision Bound]
\label{cor:T2_specialization}
Under the binary-observation model of v8 (y_t $\in$ \{0,1\}, with
$\sigma(y=1 \mid \gamma) = (1+\gamma)/2$), the per-sample Fisher information
for $\gamma$ is $\mathcal{I}_{11}/N = 1 / [4\sigma(1-\sigma)]$, giving
\[
  \mathrm{Var}(\hat\gamma) \;\ge\; \frac{4\sigma(1-\sigma)}{N},
  \qquad
  \mathrm{CV}(\hat\gamma) := \frac{\sqrt{\mathrm{Var}(\hat\gamma)}}{\gamma}
    \;\ge\; \frac{2\sqrt{\sigma(1-\sigma)}}{\gamma \sqrt N}.
\]
This is the v8 form $\mathrm{CV} \ge 1/(\gamma \sqrt{N \sigma(1-\sigma)})$
(note the factor-of-2 difference stems from v8's slightly different bound
convention~\citep[Eq.~(3.11)]{lehmanncasella1998}).
\end{corollary}

\begin{theorem}[Fisher Information Geometry: Curvature Bound]
\label{thm:T2_geometry}
Under the conditions of Theorem~\ref{thm:T2}, equip the parameter manifold
$\Theta \cong \mathbb{R}^2$ with the \textbf{Riemannian metric}
$g_{ij}(\theta) := \mathcal{I}_{ij}(\theta)$ (Amari's
information-geometric metric~\citep{amari1985}). Let $\ell(\theta)$ be the
log-likelihood of N observations. The \textbf{expected Hessian}
$-\mathbb{E}[\nabla^2_\theta \ell(\theta)]$ equals $\mathcal{I}(\theta)$.
\begin{enumerate}[nosep,label=(\alph*)]
  \item The squared geodesic distance from $\theta_0$ to $\theta$ satisfies
        \[
          d_g^2(\theta_0, \theta)
            \;\le\; \frac{1}{N} \mathbb{E}_\theta\!\left[
              \log \frac{p(y \mid \theta)}{p(y \mid \theta_0)}
            \right]^2.
        \]
  \item The curvature scalar $R(\theta)$ of $(\Theta, g)$ governs the rate at
        which $\hat\theta$ concentrates: asymptotically,
        $\sqrt{N}\,(\hat\theta - \theta) \Rightarrow \mathcal{N}(0, \mathcal{I}(\theta)^{-1})$,
        and the curvature bounds the rate of convergence of the empirical
        distribution.
\end{enumerate}
\end{theorem}

\begin{proof}
The information-geometric interpretation of the Fisher metric is classical
(Amari 1985, Amari 2016). The squared-geodesic bound in (a) follows from
Pinsker's inequality applied along the unique geodesic from $\theta_0$ to
$\theta$:
\[
  d_g^2(\theta_0, \theta)
    = \int_0^1 g(\gamma(t)) \dot\gamma(t)^2 dt
    \le \int_0^1 \mathbb{E}[(\partial_t \log p)^2] dt
    \le \frac{1}{N}\bigl(\mathbb{E}[\log(p(y\mid\theta)/p(y\mid\theta_0))]\bigr)^2,
\]
where the last step uses Jensen's inequality on the convex function
$\log^2$ and the Donsker--Varadhan representation of KL divergence.
Statement (b) is the standard asymptotic normality of the MLE on regular
families~\citep{lehmanncasella1998}, Chapter~6. \qedhere
\end{proof}

\begin{remark}[Empirical Falsification Condition]
\label{rem:T2_falsification}
Theorem~\ref{thm:T2} is refuted if there exists an unbiased estimator
$\hat\theta$ on the 11-condition registry with covariance matrix violating the
matrix inequality $\Sigma_{\hat\theta} \succeq \mathcal{I}^{-1}$. Numerically,
the bound is checked by comparing the empirical covariance of
$(\hat\gamma, \hat H)$ across the $N=30$ seeds against the closed-form
$\mathcal{I}^{-1}$ from the score function of the assumed observation model.
Pre-registered refutation: any $(i,j)$ entry violating $\Sigma \succeq
\mathcal{I}^{-1}$ by more than a one-standard-deviation sampling fluctuation.
\end{remark}

\begin{remark}[Operational Implication for v8's CV Bound]
The v8 bound $\mathrm{CV} \ge 1/(\gamma \sqrt{N \sigma(1-\sigma)})$ is
recoverable as the diagonal of Theorem~\ref{thm:T2}'s matrix inequality
(Corollary~\ref{cor:T2_specialization}). The off-diagonal entry
$[\mathcal{I}^{-1}]_{12}$ is novel: it bounds how rapidly a decrease in
$\mathrm{Var}(\hat\gamma)$ forces an increase in $\mathrm{Var}(\hat H)$,
capturing the joint geometry of the two coupled parameters. In the
information-geometric view (Theorem~\ref{thm:T2_geometry}), this coupling is
the curvature of the $(\gamma, H)$ manifold under the Fisher metric, and
provides a coordinate-free reason why CV reductions alone are insufficient
to break the impossibility triangle.
\end{remark}

% =============================================================================
% Cross-references and acknowledgments
% =============================================================================

% --- Bibliography entries used in this proof ---
%
% @book{lehmanncasella1998,
%   author = {Lehmann, Erich L. and Casella, George},
%   title = {Theory of Point Estimation},
%   year = {1998},
%   edition = {2nd},
%   publisher = {Springer}
% }
%
% @article{amari1985,
%   author = {Amari, Shun-ichi},
%   title = {Differential Geometrical Methods in Statistics},
%   year = {1985},
%   journal = {Lecture Notes in Statistics},
%   volume = {28},
%   publisher = {Springer}
% }
%
% @book{amari2016,
%   author = {Amari, Shun-ichi},
%   title = {Information Geometry and Its Applications},
%   year = {2016},
%   publisher = {Springer}
% }